\begin{afterword}
\summaryhead

The post argues that 0 and 1 should not count as probabilities, just as infinity does not count as an integer. Probabilities can be rewritten as odds, $O = P/(1-P)$, without loss, and odds make a Bayesian update a matter of multiplying by the likelihood ratio. Cox's theorem, the post says, shows that all reasonable ways of representing uncertainty are equivalent. On the odds scale, probability 1 becomes infinity.

On log odds, measured in decibels as Jaynes recommended, 0 and 1 become minus and plus infinity, and the distance between two degrees of belief equals the evidence needed to move from one to the other. So certainty would take infinitely strong evidence. Theorems of probability also need special cases at 0 and 1, such as updating on an observation of probability 0. The post proposes dropping 0 and 1, grants that theorems relying on probabilities that sum to 1 would have to be rederived, and hopes to do this without a ``magic symbol'' for all the possibilities one has not considered.

\responsehead

The post's central fact is correct and useful. In log odds, Bayes's rule is addition, so the gap between two degrees of belief is the evidence needed to cross it. It follows that from a belief short of certainty, no evidence that the alternative leaves possible can produce certainty. The worked examples are right to the digits given.

The title claims more than this. The argument shows that certainty cannot be reached by updating. It does not show that no proposition may have probability 1. The argument from scales does not supply the missing step. On odds, 1 looks remote but 0 still goes to 0, so what looks reachable depends on the scale chosen. Cox's theorem, which the post cites, puts certainty on the scale: in Jaynes's presentation, ``Certainty is represented by w(A|C) = 1.'' A change of scale changes how certainty is written, not whether it has a place.

The proposal also conflicts with the foundations it would revise. That the probabilities of all possibilities sum to 1 is not an assumption used in a few proofs. It is one of Kolmogorov's axioms, and it gives probability 1 to $A$ or not-$A$ in every model. The ``magical symbol'' the post wants to avoid is the sample space of standard theory, an ordinary event with probability 1. The analogy with infinity fails at this point: infinity is kept out of the real numbers because it breaks their axioms, while the axioms of probability require 0 and 1.

Probability 0 is needed in ordinary cases too. For a quantity with a continuous distribution, every exact value has probability 0, yet some value occurs, and standard theory updates on such observations with densities, with known pitfalls. Nor could 0 be excluded there: at most countably many exclusive events can have positive probability. In standard theory, probability 0 does not mean impossible, and probability 1 does not mean certain. The post treats them as ``infinite improbability'' and ``infinite certainty'' without drawing that distinction.

The post ends with a wish rather than a construction. It proposes no axiom to replace the one it would drop and derives no theorem without it.

In short: a correct and useful point, that finite evidence cannot carry a belief to certainty, stretched into a title claim that would mean giving up an axiom of probability, which the post admits and for which it offers no replacement.
\end{afterword}
